\originalchapter{本科期末考试: 独立解答}
试卷日期2012年12月19日, 五题共100分. 官方 OCW 页面未提供这份期末的解答, 本章全部解答为\textbf{AI 独立解答}, 经独立复核后仍不代表 MIT 官方解答. 题中给定的场和公式可在原试卷及独立的 Final Formula Sheet 核对.

\originalsection{旋转带电球壳}
\begin{exercise}
题1 (20分). 球壳半径 $R$, 均匀带电 $Q$, 绕 $z$ 轴以角速度 $\omega$ 旋转. 给定内场 $\bb_\mathrm{in}=b_0\omega\hat{\vct z}$, 外场 $\bb_\mathrm{out}=b_0\omega R^3(2\cos\theta\hat{\rr}+\sin\theta\hat{\vct\theta})/(2r^3)$, $b_0=\mu_0Q/(6\pi R)$. (a, 7分) $\dot\omega\ne0$ 时求表面感应电场与力矩; (b, 6分) 恒角速度时求磁能, 证明内外占比为2:1, 写成 $I_\mathrm{mag}\omega^2/2$; (c, 7分) 求场角动量并核对 $I_\mathrm{mag}\omega$.
\end{exercise}
\begin{solution}
(a) 所给磁场是准静态场, $\dot\omega$ 变化足够慢时按此求感应场. 用纬线包围的水平圆盘, 磁通 $\Phi=\pi R^2\sin^2\theta\,b_0\omega$. 轴对称使环路电场恒为 $E_\phi$, 法拉第定律给 $2\pi R\sin\theta E_\phi=-\pi R^2\sin^2\theta b_0\dot\omega$, 故 $E_\phi=-b_0R\dot\omega\sin\theta/2$. 面密度 $\sigma=Q/(4\pi R^2)$, 轴向力矩
\[
\tau_z=\int R\sin\theta\,\sigma E_\phi R^2\dd\Omega
=-\frac{Qb_0R^2}3\dot\omega=-\frac{\mu_0Q^2R}{18\pi}\dot\omega.
\]
(b) 记 $B_0=b_0\omega$. 内部能量 $U_\mathrm{in}=2\pi B_0^2R^3/(3\mu_0)$. 外场平方为 $B_0^2R^6(1+3\cos^2\theta)/(4r^6)$, 角积分为 $8\pi$, 径向积分为 $\int_R^\infty r^{-4}\dd r=1/(3R^3)$, 故 $U_\mathrm{out}=\pi B_0^2R^3/(3\mu_0)$. 总能 $U=\pi B_0^2R^3/\mu_0=\mu_0Q^2R\omega^2/(36\pi)$, 因而 $I_\mathrm{mag}=\mu_0Q^2R/(18\pi)$.
(c) 静电场球内零、球外径向. 第8章的积分给 $L_z=QB_0R^2/3=I_\mathrm{mag}\omega$, 横向分量由轴对称性为零. 若需用本式解释变角速度, 辐射损失不属于该准静态惯量模型.
\end{solution}

\originalsection{磁单极模型的动量}
\begin{exercise}
题2 (20分). 仅有磁荷 $\rho_m$ 和磁电流 $\jj_m$, 没有电荷与电流. 场满足 $\nabla\cdot\bb=\mu_0\rho_m$, $\nabla\times\ee=-\mu_0\jj_m-\partial_t\bb$, $\nabla\cdot\ee=0$, $\nabla\times\bb=c^{-2}\partial_t\ee$. 磁荷力为 $q_m(\bb-\vct v\times\ee/c^2)$. (a, 5分) 写力密度; (b, 5分) 用场消去源; (c, 10分) 证明机械动量加场动量的变化可写应力面积分, 是否要改应力张量?
\end{exercise}
\begin{solution}
(a) $\vct f=\rho_m\bb-c^{-2}\jj_m\times\ee$. (b) 从法拉第式求 $\jj_m=-(\nabla\times\ee+\partial_t\bb)/\mu_0$, 得
\[
\vct f=\frac1{\mu_0}\bb(\nabla\cdot\bb)
+\eps[(\nabla\times\ee)\times\ee+(\partial_t\bb)\times\ee].
\]
加上 $\partial_t\vct g$, $\vct g=\eps\ee\times\bb$, 两个 $\partial_t\bb$ 叉积抵消, 余下 $\eps\partial_t\ee\times\bb=(\nabla\times\bb)\times\bb/\mu_0$.
(c) 用 $(\nabla\times\vct F)\times\vct F=(\vct F\cdot\nabla)\vct F-\frac12\nabla F^2$, 并用 $\nabla\cdot\ee=0$, 得
\[
f_i+\partial_tg_i=\partial_j\left[\eps(E_iE_j-\tfrac12\delta_{ij}E^2)
+\mu_0^{-1}(B_iB_j-\tfrac12\delta_{ij}B^2)\right].
\]
所以应力张量无需改动. 对包围相同物质系统、无粒子跨越边界且其他外力另计的区域, 体积分成为 $\dd(\vct P_\mathrm{mech}+\vct P_\mathrm{field})/\dd t=\oint\vct T\cdot\nn\dd a$. 对一般固定控制体, 先保留 $\int\vct f$ 并另计机械动量通量.
\end{solution}

\originalsection{匀速点电荷}
\begin{exercise}
题3 (15分). 点电荷 $q$ 以亚光速 $v_0$ 沿 $x$ 轴运动, $\rr_p(t)=v_0t\hat{\vct x}$. 观察者在 $(0,y_0,0)$, $y_0\ne0$, 于 $t=0$ 测量. (a, 5分) 求推迟时间、$\rr-\rr_p$ 和 $\vct u=c\nn-\vct v$; (b, 5分) 从点电荷推迟场式求电场, 不能只引用已知匀速结论; (c, 5分) 求磁场.
\end{exercise}
\begin{solution}
设 $Y=|y_0|$, $s_y=y_0/Y$, $\beta=v_0/c$, $\gamma=(1-\beta^2)^{-1/2}$. 由 $c^2t_r^2=v_0^2t_r^2+y_0^2$ 取负根, 得 $t_r=-\gamma Y/c$. 所以 $\vct R=\gamma\beta Y\hat{\vct x}+y_0\hat{\vct y}$, $R=\gamma Y$, $\nn=\beta\hat{\vct x}+s_y\hat{\vct y}/\gamma$, $\vct u=cs_y\hat{\vct y}/\gamma$.

加速度为零, 第7章速度场给 $\kappa=1-\nn\cdot\vct\beta=\gamma^{-2}$ 与 $\nn-\vct\beta=s_y\hat{\vct y}/\gamma$. 逐项代入,
\[
\ee=\frac q{4\pi\eps}\frac{\gamma^{-2}(s_y/\gamma)\hat{\vct y}}
{\gamma^{-6}\gamma^2Y^2}
=\frac{q\gamma y_0}{4\pi\eps|y_0|^3}\hat{\vct y}.
\]
磁场为 $\nn\times\ee/c$, 等于 $v_0E_y\hat{\vct z}/c^2$. 电场指向电荷当前横向位置差, 与推迟位置矢量 $\vct R$ 的方向不同, 不能把 $\nn$ 直接当电场方向.
\end{solution}

\originalsection{旋转磁偶极}
\begin{exercise}
题4 (20分). 理想磁偶极 $\mm=m_0[\cos\omega t\hat{\vct x}+\sin\omega t\hat{\vct y}]$. (a, 10分) 求 $(0,0,z_0)$ 处的时变电、磁场; (b, 5分) 求该点坡印廷矢量及其时间平均; (c, 5分) 求 $(x_0,0,0)$ 处场. 观察点与原点不同, 原题没有指定仅在辐射区.
\end{exercise}
\begin{solution}
保留第10章 $r^{-3},r^{-2},r^{-1}$ 三阶磁偶极场, 不擅自把题目改成远场题. 记 $K=\mu_0/(4\pi)$, 对每个观察点用 $R$ 为距原点距离、$t_r=t-R/c$, $C=\cos\omega t_r$, $S=\sin\omega t_r$.

(a) 对 $z$ 轴点取 $R=|z_0|$, $s_z=\operatorname{sgn}z_0$, 定义 $\vct e_r=C\hat{\vct x}+S\hat{\vct y}$, $\vct e_t=-S\hat{\vct x}+C\hat{\vct y}$. 由于 $\nn\cdot\mm=0$, 三阶公式给
\begin{align*}
\bb&=Km_0\left[(\omega^2/(c^2R)-R^{-3})\vct e_r
-\omega\vct e_t/(cR^2)\right],\\
\ee&=-s_zKm_0\left[\omega\vct e_r/R^2+\omega^2\vct e_t/(cR)\right].
\end{align*}
(b) $\vct e_r\times\vct e_t=\hat{\vct z}$. 叉积中的 $R^{-4}$ 两项抵消, 留下
\[
\vct S=\frac{\mu_0m_0^2\omega^4}{16\pi^2c^3R^2}s_z\hat{\vct z}.
\]
它不随时间变化, 所以时间平均与自身相同. 若只在辐射区取场, 也得到此能流, 但(a)的完整场含近场与感应场.

(c) 对 $x$ 轴点取 $R=|x_0|$, $s_x=\operatorname{sgn}x_0$. 代入 $\nn=s_x\hat{\vct x}$, 得
\begin{align*}
\bb=Km_0\{&[2C/R^3-2\omega S/(cR^2)]\hat{\vct x}\\
&+[-S/R^3-\omega C/(cR^2)+\omega^2S/(c^2R)]\hat{\vct y}\},\\
\ee={}&s_xKm_0[\omega C/R^2-\omega^2S/(cR)]\hat{\vct z}.
\end{align*}
在 $\omega\to0$ 时电场消失, 磁场恢复静态偶极场; 这是对三阶表达式的符号核对.
\end{solution}

\originalsection{规范、四极与球对称辐射}
\begin{exercise}
题5 (25分). (a, 5分) 任给标量函数 $V_\mathrm{Alan}$, 能否选规范使 $V'=V_\mathrm{Alan}$, 给出 $\Lambda$? (b, 5分) 任给矢量函数 $\vct A_\mathrm{Alan}$, 能否总使 $\vct A'=\vct A_\mathrm{Alan}$? (c, 7分) 理想静态电四极的势 $V=Q_{ij}n_in_j/(8\pi\eps r^3)$, $Q$ 无迹对称, 求分布电荷密度. (d, 8分) 均匀带电球壳总电荷 $Q$, 半径 $R(t)=R_0+R_1\sin\omega t$, $0<R_1<R_0$, 证明不辐射.
\end{exercise}
\begin{solution}
(a) 在没有另加时空边界限制的局部规范问题中可以. 取 $\Lambda(\rr,t)=\int_{t_0}^t[V(\rr,s)-V_\mathrm{Alan}(\rr,s)]\dd s+\Lambda_0(\rr)$, 则 $V-\partial_t\Lambda=V_\mathrm{Alan}$. 任意 $\Lambda_0$ 仍是自由度. 若另外要求规范函数在两端时间都为零, 则会增加相容条件, 这不是原题所设限制.
(b) 不能总做到, 因为规范不改变 $\bb=\nabla\times\vct A$. 至少必须有 $\nabla\times\vct A_\mathrm{Alan}=\bb$. 单连通区域中, 差场无旋是存在标量梯度的充分条件; 非单连通域还要核对环路积分. 任给矢量函数一般不满足这些条件.
(c) 因 $Q_{ii}=0$, 接触项与 $\delta_{ij}$ 收缩消失, 第4章恒等式给 $V=Q_{ij}\partial_i\partial_j(1/r)/(24\pi\eps)$. 泊松方程及 $\nabla^2(1/r)=-4\pi\delta^3$ 给
$\rho=Q_{ij}\partial_i\partial_j\delta^3(\rr)/6$. 系数依本稿四极矩定义, 不能省去 $1/6$.
(d) 令 $\Theta$ 为阶跃函数, 精确球对称解为
\[
\ee=\frac Q{4\pi\eps r^2}\Theta(r-R(t))\hat{\rr},\qquad \bb=0.
\]
电荷与电流分别为 $\rho=Q\delta(r-R)/(4\pi R^2)$, $\jj=\rho\dot R\hat{\rr}$. 高斯定律给该面电荷, $\nabla\times\ee=0$ 满足法拉第定律. 又 $\eps\partial_t\ee=-Q\dot R\delta(r-R)\hat{\rr}/(4\pi r^2)=-\jj$, 所以安培--麦克斯韦方程也成立. 外部电场恒定、磁场为零, $\vct S=0$, 不向无穷远输出能量. 加速电荷的辐射可以在整个球对称连续分布中相消, 不能逐个加上各点电荷的辐射功率.
\end{solution}

实际带电物质还须满足 $|\dot R|<c$, 例如原题正弦脉动需 $\omega R_1<c$; 这是物理模型的补充限制, 并非原试卷明示条件.
